\section{Skills：可复用知识的渐进式披露}

\subsection{Skills 的概念}

Skills 最初由 Anthropic 为 Claude Code 引入，现已成为多个 harness（如 Codex、OpenCode）支持的开放标准。每个 skill 是一个目录，包含一个 \texttt{SKILL.md} 文件和可能的附属资源。

\begin{warningbox}{Skills 安全风险}
Skill 注册中心已经被发现分发过数百个恶意 skill。像对待 \texttt{npm install random-package} 一样对待 skill ---阅读你正在安装的内容。ClawHub、skills.sh 等注册中心的 skill 可以在你的机器上执行任意代码。
\end{warningbox}

\subsection{渐进式披露的核心思想}

HumanLayer 早期犯过的错误：把所有指令和工具都塞进 system prompt，结果 agent 越来越差。在 agent 开始工作之前，instruction budget 就已经被消耗殆尽。

\textbf{Skills 通过渐进式披露解决这个问题}：agent 只在决定需要时（或你替它决定时）才获得特定的指令、知识或工具。

\subsection{Skill 的激活机制}

当 skill 被激活时，其目录中的 \texttt{SKILL.md} 文件作为 user message 加载到 agent 的 context window 中。\texttt{SKILL.md} 文件可以告知 agent 该 skill 还捆绑了什么其他资源。

\begin{table}[H]
\centering
\caption{Skill 目录结构示例}
\begin{tabular}{@{}ll@{}}
\toprule
\textbf{文件/目录} & \textbf{用途} \\
\midrule
\texttt{SKILL.md} & 主指令文件，描述 skill 的功能和使用方式 \\
\texttt{response\_template.md} & 响应模板，规范 agent 的输出格式 \\
\texttt{CLIs/linear-cli} & 捆绑的命令行工具 \\
\texttt{CLIs/tunnel-cli} & 捆绑的命令行工具 \\
\bottomrule
\end{tabular}
\end{table}

更高级的渐进式披露：你可以在 skill 中捆绑多个 markdown 文件，每个文件包含不同功能或用途的信息，主 \texttt{SKILL.md} 文件告诉 agent 其他文件是什么以及何时应该读取它们。

\subsection{Skills 与工具分发}

目前无法直接在 skill 中捆绑 MCP server 或自定义 agent 工具。替代方案是将工具编写为可执行文件、CLI、NPM 包或其他可以与 skill 一起分发或在 skill 文件中指示 agent 安装的形式。

例如，与其配置一个 Playwright MCP server，不如提供一个基于 BrowserBase 的 agent browser skill 或 Vercel 的 agent browser CLI 的 web 浏览 skill。

\subsection{本章小结}

Skills 是实现渐进式披露的关键机制---只在需要时加载知识和工具，避免 system prompt 过载。使用时需注意安全风险，避免盲目安装未审核的第三方 skill。


